\Needspace{9\baselineskip}\begin{center}\textit{Divisions of the circle that can be carried out by quadratic equations or by\\geometric constructions.}\end{center}


\Needspace{6\baselineskip}\Article{365}
Through the preceding investigations we have therefore reduced the division of the circle into \(n\) parts, when \(n\) is prime, to the solution of as many equations as there are factors into which the number \(n-1\) can be resolved; the degrees of these equations are determined by the sizes of the factors. Thus, whenever \(n-1\) is a power of the number 2, as happens for \(n\) at the values 3, 5, 17, 257, 65537, etc., the division of the circle reduces to quadratic equations alone, and the trigonometric functions of the angles \(\frac Pn,\allowbreak \frac{2P}n\), etc., can be expressed by square roots of greater or lesser complexity (according to the size of \(n\)). In these cases, therefore, the division of the circle into \(n\) parts, or the drawing of a regular polygon of \(n\) sides, can evidently be accomplished by geometric con\SourceLeaf{467}structions. Thus, \(\text{e.g.}\) for \(n=17\), Articles 354 and 361 readily give the following expression for the cosine of the angle \(\tfrac1{17}P\): \[\begin{aligned}&-\tfrac1{16}+\tfrac1{16}\surd17+\tfrac1{16}\surd(34-2\surd17)\\&\quad+\tfrac18\surd(17+3\surd17-\surd(34-2\surd17)\\&\qquad\qquad-2\surd(34+2\surd17))\end{aligned}\] The cosines of multiples of this angle have a similar form, but the sines have one more radical sign. It is indeed greatly to be wondered at that, although the geometric divisibility of the circle into three and five parts was already known in Euclid's time, nothing was added to these discoveries over an interval of 2000 years, and all geometers pronounced it certain that, apart from these divisions and those that follow from them of themselves, namely divisions into \(15,\allowbreak 3.2^\mu,\allowbreak 5.2^\mu,\allowbreak 15.2^\mu\) and also into \(2^\mu\) parts, no others could be accomplished by geometric constructions. — Further, it is easy to prove that if the prime \(n\) is \(=2^m+1\), the exponent \(m\) too can contain no prime factors other than the number 2, and must therefore be either \(=1\), or \(=2\), or equal to a higher power of the number 2. For if \(m\) were divisible by any odd number \(\zeta\) (greater than unity), and \(m=\zeta\eta\), then \(2^m+1\) would be divisible by \(2^\eta+1\), and therefore necessarily composite. Thus all values of \(n\) for which we arrive at quadratic equations alone are contained in the form \(2^{2^\nu}+1\); the five numbers 3, 5, 17, 257, 65537 arise by setting \(\nu=0,\allowbreak 1,\allowbreak 2,\allowbreak 3,\allowbreak 4\), or \(m=1,\allowbreak 2,\allowbreak 4,\allowbreak 8,\allowbreak 16\). The division of the circle can by no means be carried out geometrically for \textit{all} numbers contained in that form, but only for those that are prime. Fermat, misled by induction, had indeed asserted that all numbers contained in that form must be prime; but the distinguished Euler was the first to notice that this rule is already false for \(\nu=5\), or \(m=32\), since the number \(2^{32}+1=4294967297\) contains the factor 641.

Whenever, however, \(n-1\) contains prime factors other than 2, we always arrive at higher equations: namely, at one or more cubic equations when 3 occurs once or more among the prime factors of \(n-1\), at equations of the fifth degree when \(n-1\) is divisible by 5, and so on; \textsc{and we can prove with complete rigor that these higher equations can in no way be avoided or reduced to lower ones,} although the limits of this work do not allow us to give this proof here. We nevertheless thought it necessary to mention this, so that no one should hope to achieve by geometric constructions still other divisions beyond those suggested by our theory, \(\text{e.g.}\) divisions into 7, 11, 13, 19, etc., parts, and waste time in vain.
\SourceLeaf{468}


\Needspace{6\baselineskip}\Article{366}
If the circle is to be divided into \(a^\alpha\) parts, where \(a\) denotes a prime, this can evidently be accomplished geometrically when \(a=2\), but for no other value of \(a\), provided \(\alpha>1\). For then, besides the equations required for division into \(a\) parts, it is necessary to solve \(\alpha-1\) further equations of degree \(a\); these too can in no way be avoided or reduced in degree. Thus the degrees of the necessary equations can be known in general from the prime factors of the number \((a-1)a^{\alpha-1}\) (that is, also in the case \(\alpha=1\)).

Finally, if the circle is to be divided into \(N=a^\alpha b^\beta c^\gamma\ldots\) parts, where \(a,\allowbreak b,\allowbreak c\), etc., denote unequal primes, it suffices to have carried out the divisions into \(a^\alpha,\allowbreak b^\beta,\allowbreak c^\gamma\), etc., parts (Article 336). To know the degrees of the equations needed for this purpose, we must therefore consider the prime factors of the numbers \[(a-1)a^{\alpha-1},\quad(b-1)b^{\beta-1},\quad(c-1)c^{\gamma-1}\ \text{etc.}\] or, what here amounts to the same thing, of their product. Observe that this product expresses the number of numbers coprime to \(N\) and smaller than it (Article 38). The division can therefore be accomplished geometrically only when this number is a power of two; when, however, it contains prime factors other than \(2\), say \(p,\allowbreak p'\), etc., equations of degree \(p,\allowbreak p'\), etc., can in no way be avoided. From this we conclude in general that, for the circle to be divisible geometrically into \(N\) parts, \(N\) must be \textit{either} \(2\) or a higher power of \(2\), \textit{or} a prime of the form \(2^m+1\), \textit{or} a product of several such primes, \textit{or} the product of one or more such primes with \(2\) or a higher power of \(2\). More briefly, it is required that \(N\) contain neither an odd prime factor that is not of the form \(2^m+1\), nor any prime factor of the form \(2^m+1\) more than once. The following values of \(N\) occur below \(300\), \(38\) in all: \(2,\allowbreak 3,\allowbreak 4,\allowbreak 5,\allowbreak 6,\allowbreak 8,\allowbreak 10,\allowbreak 12,\allowbreak 15,\allowbreak 16,\allowbreak 17,\allowbreak 20,\allowbreak 24,\allowbreak 30,\allowbreak 32,\allowbreak 34,\allowbreak 40,\allowbreak 48,\allowbreak 51,\allowbreak 60,\allowbreak 64,\allowbreak 68,\allowbreak 80,\allowbreak 85,\allowbreak 96,\allowbreak 102,\allowbreak 120,\allowbreak 128,\allowbreak 136,\allowbreak 160,\allowbreak 170,\allowbreak 192,\allowbreak 204,\allowbreak 240,\allowbreak 255,\allowbreak 256,\allowbreak 257,\allowbreak 272.\)
\SourceLeaf{469}


